\section{Motivation and setup}

Feature scaling changes the relative penalties in regularized regression.
\citet{larsson2025normalization} study this effect for binary features and
obtain explicit normalization results under orthogonality. The present
investigation asks which conclusions extend to dependent features and which
conditions can replace orthogonality. Our initial focus is on the distinction
between Pearson correlation and covariance relative to a feature's variance.
That distinction becomes consequential when one feature is rare.

Finite-sample distributions for the lasso with correlated designs are already
available \citep{ewald2020distribution}. Polyhedral selection events are also
central to selective inference \citep{lee2016exact}. Accordingly, the candidate
contribution is a characterization of normalization for binary features,
rather than a general distribution theory for correlated lasso estimators.
The references here are a starting point for a targeted novelty review.

Profile out an unpenalized intercept and let $X_c$ and $y_c$ denote the centered
design and response. Write
\[
 C=X_c^\top X_c/n,\qquad z=X_c^\top y_c/n,\qquad
 y_c=X_c\beta^*+\varepsilon_c.
\]
Unless stated otherwise, $C$ is positive definite and inference is conditional
on a fixed design. Independent Gaussian errors with variance
$\sigma_\varepsilon^2>0$ give
\begin{equation}\label{eq:score-law}
 z\sim N(C\beta^*,\tau^2C),\qquad \tau^2=\sigma_\varepsilon^2/n.
\end{equation}
For fixed positive vectors $u,v$, consider the original-scale estimator
\begin{equation}\label{eq:objective}
 \hat\beta=\arg\min_{\beta\in\R^p}
 \left\{\frac12\beta^\top C\beta-z^\top\beta
 +t\sum_j u_j|\beta_j|+\frac r2\sum_j v_j\beta_j^2\right\},
 \qquad t,r\ge0.
\end{equation}
Here $t=\lambda_1/n$ and $r=\lambda_2/n$ correspond to penalties with an
unaveraged squared-error loss. Dividing centered feature $j$ by $s_j$ and
transforming the fitted coefficient back gives $u_j=s_j$ and $v_j=s_j^2$.
This equivalence requires no orthogonality. Allowing separate weights also
covers penalty schemes such as $u_j=v_j=C_{jj}$.

Weights and penalties are fixed with respect to $y$ throughout the distribution
calculations. Cross-validation and response-dependent penalty paths require
separate analysis. Write $\soft_k(x)=\sign(x)(|x|-k)_+$ for
soft-thresholding.

\begin{proposition}[Ridge in arbitrary dimension]\label{prop:ridge}
Let $V=\diag(v)$ and $t=0$. Then
\begin{align*}
 \E\hat\beta&=(C+rV)^{-1}C\beta^*,\\
 \E\hat\beta-\beta^*&=-r(C+rV)^{-1}V\beta^*,\\
 \Cov(\hat\beta)&=\tau^2(C+rV)^{-1}C(C+rV)^{-1}.
\end{align*}
The mean and covariance formulas require only zero-mean errors with covariance
$\sigma_\varepsilon^2 I$, rather than Gaussian errors.
\end{proposition}
\begin{proof}
Differentiate the quadratic objective to obtain $\hat\beta=(C+rV)^{-1}z$,
then take its mean and covariance.
\end{proof}
For $r>0$ and $0<k<1$, the identity $(C+rV)^{-1}C=kI$ implies
$V=(1-k)C/(kr)$. Thus a diagonal ridge penalty can give the same shrinkage
factor for every coefficient vector only when $C$ is diagonal.
